%=== True False ===

\question\TFQuestion{F}{The linear arrangement the toolkit genes \textit{Shh} and \textit{BMP4} corresponds to their proximal to distal use in the developing appendages of vertebrates.}

\question[5]
Discuss the source and role of crystallin proteins in the evolution of the vertebrate eye.
	\begin{solution}
	Based on page 315 and question 7 (pg 325), plus lecture.\\
	Crystallin proteins were co-opted from heat shock proteins. The co-opted proteins had the correct optical properties for bending light so that it could be focused.\\
	Consider grading this lightly.
	\end{solution}
\vspace*{\stretch{1}}

%%%%%
\question[10]
Whales are most closely related to hippopotamuses but lack rear limbs.  Use your knowledge of \emph{hox} genes to hypothesize a genetic mechanism that could explain how the rear limbs were lost. Be as precise as possible with the specific set(s) of possible \emph{hox} genes.
	\begin{solution}
	Hox A and Hox D 9-13 may have stopped limb production. Alternatively, hox genes for the a-p axis may have stopped producing legs in the posterior end.
	\end{solution}

\vspace*{\stretch{1}}

%%%%
\question[8]
The anglerfish (pictured on screen) has an unusually wide mouth compared to other species of anglerfishes. You wish to study the genetics of jaw morphology in the anglerfish.  
\begin{parts}

\part[2] Identify one toolkit gene that you would be wise to study first.


\vspace*{\stretch{0.5}}

\part[5] Explain the existing evidence that caused you to identify this toolkit gene (5 pts).
\vspace*{\stretch{1}}
\end{parts}

%% Do one for duck-billed platypus: Wide, flat.

\question[7]
The longnose butterflyfish (right, on screen) has an unusually long mouth compared to other species of butterflyfishes. You plan to study the genetics of jaw length in the butterflyfishes.  
\begin{parts}
	\part Identify one toolkit gene that you would be wise to study first (2 pts).
	\begin{solution}
		CaM1
	\end{solution}
	\vspace*{\stretch{0.5}}
	
	\part Explain the existing evidence that caused you to identify this toolkit gene (5 pts).
	\begin{solution}
		CaM1 (expression) was shown to correlate with jaw width of cichlid fishes.
	\end{solution}
	\vspace*{\stretch{1}}
\end{parts}





\question[5]
The so-called ``Cambrian Explosion'' was named in the 1800s because the known fossil record of that time suggested that all of the major phyla appeared almost entirely at the start of the Cambrian. Recent and detailed fossil evidence, coupled with genetic evidence, suggests that ``Cambrian Radiation'' is a better term.  Describe one type of fossil or genetic evidence that explains why evolution across the Ediacaran-Cambrian boundary was not ``explosive.'' 
	\begin{solution}
	Molecular clock show most phyla evolved before Cambrian.\\
	Fossils extending back to Ediacaran\\
	Fossils accumulating over 30 million year period at start of Cambrian.
	\end{solution}	
\vspace*{\stretch{1.5}}

\newpage

\question[5]
Why was cyanobacteria photosynthesis was important to the evolution of increased body size for single-celled eukaryotes?
	\begin{solution}
	Increased oxygen concentration in the atmosphere, sufficient for cellular respiration in larger organisms.
	\end{solution}	
\vspace*{\stretch{0.75}}

\question[10]
The snapping shrimp \textit{Synalpheus brooksi} lives in sponges on coral reefs.  Duffy (1996) showed that shrimp occupying nearby but different species of sponges rarely come into contact with each other because they don't usually leave their sponges. In the lab, a male shrimp from one sponge species is extremely aggressive towards females from a different sponge species, suggesting reproductive incompatibility. Analyses showed that shrimp from the same sponge species were genetically similar but shrimp from different sponge species were genetically distinct. Duffy argued the shrimp in different sponges are different species.  If so, explain which prezygotic mechanism or mechanisms reproductively isolate shrimps in different sponge species.
	\begin{solution}
	Ecological. Occupy different habitats.
	\end{solution}

\vspace*{\stretch{1.5}}

\newpage

\question[10]
The grasses (Poaceae, includes rice and corn) produce flowers that lack petals and sepals but has stamens and carpels.  Given the three classes of \textsc{mads} genes below, which class or classes most likely had a mutation or change in expression that caused the loss of petals and sepals?  Draw the ABC model of floral develop, and then use that model to justify your answer.  Class A: Apetala1.  Class B: Apetala3.  Class C: Agamous.
	\begin{solution}
	Class A causes loss of sepals and petals. Class B and C form stamens and C alone forms carpals. A needed for both petals and sepals. 
	\end{solution}
	
\vspace*{\stretch{1}}

\question[10]
Whales are most closely related to hippopotamuses but lack rear limbs.  Explain how a change to \emph{hox} gene expression might cause the loss of rear limbs. Identify the specific \emph{hox} genes involved, including the specific set of paralogs (duplicated copies of the genes) and numbers (e.g., Hox-A1, Hox-B13, Hox-C2-5, etc.).

	\begin{solution}
	Mutations to Hox A and Hox D 9-13 may have stopped limb production. Alternatively, hox genes for the a-p axis may have stopped producing legs in the posterior end.
	\end{solution}

\vspace*{\stretch{1}}


\question[5]
Whales are most closely related to hippopotamuses but lack rear limbs.  Explain how a change to \emph{Box} gene expression might cause the loss of rear limbs. 

\begin{solution}
	Mutations to Hox A and Hox D 9-13 may have stopped limb production. Alternatively, hox genes for the a-p axis may have stopped producing legs in the posterior end.
\end{solution}

\vspace*{\stretch{1}}




\fullwidth{The phylogeny below shows the relationships of gobies (fishes) in the genus \textit{Elacatinus}. Each species has a stripe on that side. The color of the stripe for each species is in parentheses. The vertical line indicates the behavior of each species (sponge-dwelling, hovering, cleaning). Every species to the left of the vertical line has that behavior. Based on this information, answer the next two questions:}
\begin{center}
\includegraphics[width=0.45\textwidth]{goby_tree}
\end{center}

\question[5]
Circle the \emph{largest} monophyletic group you can identify based on color \emph{or} behavior. Be clear which species are included / excluded. Explain why that group is monophyletic. Do not circle the entire tree. That would be incorrect.
\vspace*{\stretch{1}}

\question[10]
Does color or behavior better explain most differences among closely related species? Explain why, providing specific examples. 
	\begin{solution}
	Color. Most sister species differ by coloration. Behavior forms two large groups and two independent evolutions of hovering. Example counts for 4 points.
	\end{solution}
\vspace*{\stretch{2}}


\newpage

\question[20]

The purple-throated carib is a species of hummingbird that lives on the Caribbean island of Dominica. Female caribs have longer, more curved bills than male caribs, which have shorter, straighter bills. Female hummingbirds get their nectar only from the red flowers of a plant called \textit{Heliconia bihai}. Male hummingbirds get their nectar only from the yellow flowers of the same plant species. An individual plant has either red or yellow flowers, but not both. The flower shape of each color morph corresponds to the bill shape of each sex. Red flowers have a longer, more curved flower corolla, while yellow flowers have a shorter, straighter corolla. The hummingbirds function as pollinators for \textit{H. bihai}.  No other pollinators (including wind) for this plant are known. Based on this information only, can reproductive isolation evolve for (1) the hummingbirds and (2) the color morphs of \textit{H. bihai}.  If so, say how. Justify your answer. 
	\begin{solution}
	Yes for H. carib because plant has different pollinators. 10 points\\
	No for birds because males and females must come together to reproduce. 10 points.
	\end{solution}



%%%